\subsection{Inferring the Objective: IRL and CIRL}
\label{sec:4.2.3}
Inverse reinforcement learning runs section~\ref{sec:4.2.2}'s method backward. Instead of taking the reward as given and learning the behavior, it takes the behavior as given and infers the reward that best explains it — watching a driver's speed choices, following distances and reactions to weather, and recovering the priorities behind them without being handed a rulebook. Cooperative inverse reinforcement learning, formalized by Dylan Hadfield-Menell and colleagues \autocite{hadfieldmenell2016cooperative}, goes one step further: the system treats the human's true objective as unknown, holds its own estimate as provisional, and keeps revising it through the interaction rather than fixing it at training time.

That last property is why CIRL is attractive, and it is also why this book cannot build a floor on it. A system that treats what it should want as permanently open, to be settled by whoever it is interacting with, is the most corrigible design in this chapter. Section~\ref{sec:3.5} argues that corrigibility and the capacity to hold a line are one property with the sign flipped — a system reliably steerable by whoever holds it is, for that reason, a system that cannot hold anything against them. CIRL is that argument in its purest engineering form. Uncertainty about the objective is exactly what makes a system defer, and deference to the party in possession is what a floor exists to interrupt.

The practical version of the problem is narrower and does not depend on chapter~\ref{sec:3} being right. Both methods learn from demonstrations, so both inherit whoever demonstrated: a narrow set of drivers, a narrow set of households, a narrow set of decisions that somebody had occasion to record. Unexamined bias in that record is learned rather than corrected, and neither method contains anything that would notice.

What survives is a real and bounded use. Above the floor, where the question is what a person actually wants rather than what may never be done to them, inferring preferences from behavior beats asking somebody to write them down, because most of what people want has never been written down and would be misstated if it were.
